\begin{abstract}
The dominant trajectory of modern machine learning has been to
scale \emph{up}: larger models, larger accelerators, larger memory
budgets. Yet a multi-year global semiconductor supply
constraint~\cite{burkacky2022chipshortage,khan2021chipshortage} and
the growing energy and carbon cost of always-online
inference~\cite{strubell2019energy,patterson2021carbon} expose the
fragility of this trajectory and motivate the opposite direction:
\emph{refactoring} AI and ML algorithms to fit the small,
ubiquitous microcontrollers that are already in mass production in
wearables, sensors, and edge appliances. This
paper is a study in that direction. We present an end-to-end
open-source reproduction of
\fastgrnn{}~\cite{kusupati2018fastgrnn}, a compact gated recurrent
cell, deployed on two bare-metal targets representative of the
low end of the available silicon supply: the 8-bit \arduino{}
(ATmega328P) and the 16-bit \msp{} (no hardware multiplier;
16~KB Flash; 512~B SRAM). Our compression pipeline combines
low-rank weight factorization, iterative hard-thresholding
sparsity, and per-tensor Q15 post-training quantization with
explicit activation calibration. The deployed model occupies
\textbf{566~bytes of weights}, achieves \textbf{macro F1 $=$ 0.918}
on the \hapt{} benchmark, and matches a PyTorch reference at
\textbf{100\% prediction agreement} across 3{,}399 test windows.
Both platforms sustain \textbf{real-time 50~Hz streaming
inference} (9.21~ms per sample on Arduino; 13~ms on MSP430), where
a 256-entry sigmoid/tanh look-up table delivers a
\textbf{30.5$\times$ speedup} on the multiplier-less MSP430.
Four contributions extend the original \fastgrnn{} paper:
(i) cross-platform bit-equivalent deterministic inference;
(ii) characterization of a $\sim$2~s recurrent warm-up latency that
bounds end-to-end response time;
(iii) a deployable look-up-table recipe for multiplier-less embedded
targets delivering a \textbf{30.5$\times$} end-to-end speedup; and
(iv) a hardware energy characterization showing \textbf{17.7~mW}
active inference power, \textbf{$<$0.09~mW} idle power, and a
\textbf{96.7\%} reduction in energy per inference window with the LUT.
Together, these results demonstrate that compact recurrent
architectures, combined with calibrated quantization, look-up-table
activations, and measured energy profiling, can deliver accurate,
memory-efficient, and energy-efficient human activity recognition on
ultra-resource-constrained microcontrollers without specialized
accelerators.
\end{abstract}
